% ============================================================
%  ch08_discussion.tex — Discussion
%  Requires: macros.tex
%  Estimated length: 18–25 pages
%  Rewritten: VD-06W, 2026-03-07
% ============================================================

\chapter{Discussion}\label{ch:discussion}

The numerical results of Chapter~\ref{ch:results} invite three
levels of interpretation.  At the mathematical level, the
master defect identity provides exact algebraic constraints on
any Bianchi cosmology; no interpretation is required.  At the
model-selection level, the Bayesian evidence comparison reveals
a clean dichotomy between CMB-only data (which prefer isotropy)
and the matter dipole (which, conditional on the anomaly being
physical, drives strong Bianchi preference); here interpretation
enters through the prior specification and the choice of data
channels.  At the physical
level, the connection between the defect framework, the Dipole
Cosmological Principle, the Tsagas tilted-observer programme,
cosmological backreaction, and the emerging observational
picture of large-scale vorticity requires a more extended
discussion that honestly delineates what is established, what is
model-dependent, and what remains speculative.

The revised interpretation is therefore deliberately narrower than a
detection narrative.  The analysis is conditional on the dipole premise:
the framework
maps which diagnostic and HTT likelihood surfaces move under current
matter-dipole inputs; it does not promote those surfaces to native
solver results, Bianchi geometry detection, or Bianchi family
identification.  The rule is explicit: family identification remains blocked
until a native low-ell morphology atlas, matched masks, covariance/null
controls, and family-equivalence tests exist.

One concrete methodological result now sits upstream of this
discussion.  Section~\ref{sec:dipole-discriminator} of
Chapter~\ref{ch:dipole} introduces an observer-frame
likelihood-ratio discriminator that keeps the local observer
boost and the cosmological tilt as separate typed surfaces, pins
their composition order, and verifies its calibration with
synthetic coverage tests.  The interpretive claims in this
chapter should therefore be read against that operational
separation: the question is no longer whether a dipole can be
written as a velocity, but whether the data prefer a local frame
correction or a genuinely cosmological tilt once the two options
are compared on the same likelihood surface.

We address these levels in order: contrastive analysis of the
seven principal results
(Section~\ref{sec:contrastive}), the DCP--defect dictionary
(Section~\ref{sec:DCP-defect}), the Tsagas programme
(Section~\ref{sec:disc-tsagas}), the backreaction caveat
(Section~\ref{sec:disc-backreaction}), cosmic-web vorticity
(Section~\ref{sec:disc-cosmic-web}), and methodological
limitations as design choices
(Section~\ref{sec:disc-limitations}).


% ============================================================
\section{Contrastive analysis of the principal results}
\label{sec:contrastive}
% ============================================================

Seven results from the preceding chapters merit separate
assessment through the conservative, progressive, and
data-honest lenses established in
Section~\ref{sec:reframing} of Chapter~\ref{ch:dipole}.
the conditioned legacy figure gallery in Appendix~\ref{app:conditioned-legacy-figure-gallery} presents these seven
results on a common FLRW departure scale.


\subsection{The CMB/non-CMB dichotomy}

The most structurally informative finding of the evidence
comparison is not the headline $\ln\mathcal{B} \approx +26.3$
but the \emph{channel decomposition}: the matter dipole
channels (b+c) alone yield $\ln\mathcal{B} \approx +29$ for
all five surviving tilted models, while the CMB-only channels
(a, d, e, f, h) return $\ln\mathcal{B} < 0$ for every model
tested.  The positive preference is therefore conditional on the
matter-dipole premise rather than an independent anisotropic-geometry
claim.

The conservative reading holds that this dichotomy exposes a
vulnerability: if the CatWISE and radio dipole excesses are
eventually attributed to unidentified survey systematics, the
legacy transfer-conditional preference for a tilt-like degree of
freedom should be retired.  The 2025
simulation-based reassessment of
CatWISE~\cite{Bashir2025}---which reduced the
CatWISE-specific significance from $4.9\sigma$ to
$\sim\!3.3$--$3.6\sigma$ by forward-modelling the WISE scanning
law and photometric systematics---illustrates precisely this
pathway.  The data do not establish anisotropy independently of
the anomaly.  The legacy transfer-conditional evidence summary is conditional on the anomaly
being substantially physical, and is one survey generation away
from confirmation, down-weighting, or refutation.

The progressive reading holds that the dichotomy is
consistent with what the Tsagas mechanism predicts: a tilt
that is invisible to the CMB (because the CMB dipole is
dominated by the kinematic component) but manifest in matter
surveys (where the tilt-induced modulation adds to the
kinematic signal).  Under this reading, the CMB's preference
for isotropy is not evidence against anisotropy but evidence
that the anisotropy resides in the tilt degree of freedom,
which the CMB constrains only weakly.  The caveat is that
the Tsagas mechanism predicts the \emph{qualitative}
pattern (CMB-invisible, matter-visible) but does not
uniquely predict the amplitude or the survey in which the
excess first appears.

The data-honest synthesis notes that the dichotomy is a
\emph{measurement}, not an interpretation: the data factually
prefer tilted geometries when matter dipole information is
included and prefer isotropy when it is excluded.  Whether the
dipole excess is physical or systematic is an empirical
question that the present framework identifies but cannot
resolve.  The conditional structure of the evidence
($\ln\mathcal{B} \approx +26.3$ if dipole is physical;
$\ln\mathcal{B} \leq 0$ if dipole is systematic) is itself a
falsifiable prediction.  Euclid DR1 (October 2026), the full
Vera Rubin LSST, and SKA Phase~1 will provide independent dipole
measurements at different redshifts and frequencies, breaking the
current reliance on a small number of surveys.


\subsection{Conditional evidence: $\ln\mathcal{B} \approx +26$ (VER05 primary)}

Six tilted models (FLRW$_{\mathrm{tilt}}$, BI, BIII, BIX,
BV, BVIIh) cluster at
$\ln\mathcal{B} \in [+24.8, +26.4]$ under the primary WFH2009
input ($\beta = 1.334 \times 10^{-3}$), corresponding to
odds exceeding $10^{11}:1$ in favour of tilted models over
FLRW, conditional on the matter-dipole anomaly being
substantially physical.  The pure-tilt
FLRW$_{\mathrm{tilt}}$ reaches $+26.3$; the Bianchi
models sit at $[+24.8, +26.4]$, with the tilt\_flat
equivalence class (BI, BIII, BIX tilt) producing
indistinguishable likelihoods.

The evidence decomposition confirms that a pure-tilt model
(FLRW$_{\mathrm{tilt}}$, with $\Sigma^2 = 0$) captures the
entire signal, while anisotropic shear alone is mildly
disfavoured.  The ``Bianchi evidence'' is therefore evidence
for a non-zero dipole anomaly/tilt, not for anisotropic
spatial geometry.

The sensitivity to the CF4 data version is substantial:
the CF4++ value (Courtois et~al.~\cite{Courtois2025})
gives $\ln\mathcal{B} \approx +44$ (sensitivity run; see
Table~\ref{tab:cf4-triple}), while
Watkins et~al.~\cite{Watkins2023} gives
$\ln\mathcal{B} \approx +106$ (smallest $\sigma_\beta$).  All three are
decisive.  The attribution error that assigned
$\beta = 1.334 \times 10^{-3}$ to Watkins+2023 (corrected
in this revision) does not affect the qualitative conclusion
but shifts the headline number by $\sim 18$ nats.

The 2025 CatWISE reassessment~\cite{Bashir2025} demonstrates
that survey-specific systematics can reduce the significance
of individual dipole measurements by $\sim\!1.5\sigma$; the
multi-survey radio combination retains $\geq 5\sigma$
independently of CatWISE~\cite{Boehme2025}.


\subsection{Bianchi~V: corrected result $\ln\mathcal{B} = +24.8$ (VER06)}

The VER04 pipeline reported BV~(tilt) at $\ln\mathcal{B} = -24.5$,
appearing to exclude the model decisively.  That result was
a code artifact: the transfer-function formula incorrectly
included a factor of $T_0 = 2.726 \times 10^6\;\mu\mathrm{K}$
when $T_2$ was already in $\mu\mathrm{K}/(\sigma/H)$ units,
compounded by a $5.24\times$ miscalibration in $T_2^{\mathrm{decay}}$.
The corrected VER06 pipeline, with fresh dynesty runs and
self-consistent $\ln\mathcal{Z}$/$\ln\mathcal{B}$ fields
(three-seed convergence verification), gives
$\ln\mathcal{B}_{\mathrm{VER06}}(\text{BV tilt}) = +24.8$
(Section~\ref{sec:BV-excl} of Chapter~\ref{ch:results}).

The conservative reading of this correction is that the framework
is honest about its own limitations: when a code error is found,
it must be corrected regardless of whether the result is
``cleaner'' before or after the fix.  The framework does
demonstrate falsifiability through BVIIh~(orth, grow) at
$\ln\mathcal{B} = -18.7$, which is a genuine Saadeh-vorticity
exclusion.

The progressive reading notes that the correction resolves a
former tension: the momentum constraint for BV~(tilt) is an
algebraic identity of the Bianchi~V structure constants,
independent of the equation of state.  For $w = 1/3$ the
prefactor $(1+w)^2$ increases by $78\%$, worsening the
overproduction.  Bianchi~V is falsified by the joint
dipole--quadrupole data, not by the dipole alone.

The data-honest caveat is that the exclusion assumes the LRS
momentum constraint holds exactly at all epochs, for a
single barotropic fluid with constant equation-of-state
parameter $w \in [0, 1/3]$, and in the absence of a
dark-energy--matter tilt coupling that could modify the
momentum balance.  In a more general (non-LRS) Bianchi~V
model, in the DCP extension with multi-fluid structure, or
with a dynamical dark-energy equation of state, the constraint
may be relaxed.  The exclusion should be reported as applying
to the LRS single-fluid Bianchi~V class.


\subsection{Filling fraction: $\FF \approx 6.3\%$}

At the S3 scenario with $w = 0$, the tilt occupies $6.3\%$ of
the MES-allowed anisotropy budget.  The Bayes factor for
$\FF > 0$ exceeds $5 \times 10^4$, which is a restatement of
the $>5\sigma$ dipole anomaly in defect language.

The conservative reading notes that $\FF \ll 1$: the universe
is overwhelmingly close to FLRW.

The progressive reading notes that $\FF$ is invariant under
frame corrections (because numerator and denominator shift by
the same factor) and provides a dimensionless measure of FLRW
departure directly comparable across scenarios.  A non-zero
filling fraction is an observational fact contingent only on
the anomaly being physical.


\subsection{Frame-attribution bias: $\etaud \approx 8.3\%$}

The acceleration dipole
$\eps_1^{(\dot{u})} = \etaud\beta$ contributes $\sim 8.3\%$
of the total dipole at S3.  The conservative reading treats
$\etaud$ as a systematic correction that tightens the
safe-route $\beta$ by $7.7\%$ and should be included in all
analyses.  The progressive reading treats $\etaud$ as a
prediction: the acceleration content of the CMB dipole
generates a specific equation-of-state-dependent relationship
between intrinsic and observed dipole amplitudes.  A measurement
of the dipole at multiple redshifts (via Euclid and SKA) would
test this prediction by measuring $\beta(z)$ and comparing with
the $\etaud$-corrected safe-route.


\subsection{Saadeh--MES gap: six orders of magnitude}

The Saadeh model-dependent bound
($\Wstd < 9.01 \times 10^{-22}$) is six orders of magnitude
tighter than the MES algebraic bound
($\Wstd < (3/2)\Bom^2 \sim 10^{-15}$), shrinking to
approximately one order for the regular tensor growing mode of
VII$_h$ (Section~\ref{sec:VIIh-modes}).

The data-honest interpretation is that the two frameworks are
complementary: MES provides type-independent algebraic ceilings
(under the almost-FLRW and Copernican assumptions)
holding for any Bianchi type, while Saadeh~et~al.\ exploit the
full morphological information of the CMB sky for specific
template shapes.  The gap quantifies the information gain from
morphological analysis.  For the regular tensor mode, where the
transfer function is $\sim 5000\times$ weaker, the information
gain collapses and the two approaches converge---a structural
consistency check.


\subsection{Preferred-direction clustering}

The anomaly directions---CMB dipole $(264^\circ, 48^\circ)$,
CatWISE $(240^\circ, -4^\circ)$, radio
$(230^\circ, -14^\circ)$, quadrupole--octupole axis
$(237^\circ, 20^\circ)$, CF4 bulk flow
$(276^\circ, 10^\circ)$---cluster within $\sim 30^\circ$ of
each other in Galactic coordinates.  The probability of five
independent directions clustering within a $30^\circ$-radius
cap by chance is $\sim 2 \times 10^{-4}$.
the conditioned legacy figure gallery in Appendix~\ref{app:conditioned-legacy-figure-gallery} shows these anomaly directions
in a Mollweide projection.

The progressive reading identifies this clustering as evidence
for a single preferred axis, consistent with the tilt direction
of a Bianchi model.  The data-honest caveat is that the
summary-statistic evidence (Chapter~\ref{ch:results}) treats
the dipole as a scalar amplitude: directional information does
not enter the seven-channel likelihood.  The preferred-axis
clustering is therefore a \emph{post-likelihood external
consistency check}, not a component of the Bayesian evidence.

The catalog pipeline
(Section~\ref{sec:catalog} of Chapter~\ref{ch:pipeline})
provides a partial directional constraint that operates at
the likelihood level.  The three-dimensional
$(\beta, l, b)$ posterior from synthetic injection-recovery
(the conditioned legacy figure gallery in Appendix~\ref{app:conditioned-legacy-figure-gallery}) recovers the injected
direction $(l, b) = (264^\circ, 48^\circ)$ within a
$68\%$ credible cone of $11^\circ$ and a $95\%$ cone of
$19^\circ$, with a posterior median at
$(240^\circ, 48^\circ)$---offset $16^\circ$ from truth.
The CMB dipole direction falls inside the $95\%$ contour.
The CatWISE direction lies near the posterior peak.  The
CF4 bulk-flow direction, which has a distinct physical
origin (gravitational infall rather than cosmological tilt),
is offset $\sim 60^\circ$ from the recovered axis.

We regard this as encouraging but not conclusive: the catalog
pipeline operates on synthetic data with known injection
parameters.  Application to real catalogs (Pantheon+,
DES-SN5YR) requires survey-specific selection functions beyond
the scope of this work.  A pixel-level Bianchi template
analysis---building on the Saadeh~et~al.~\cite{Saadeh2016b}
framework but adding the matter-dipole constraint---would
constitute the decisive directional test.


% ============================================================
\section{The Dipole Cosmological Principle in defect language}
\label{sec:DCP-defect}
% ============================================================

Our goal in this section is to construct the explicit algebraic
dictionary between the DCP dynamical variables and the defect
variables, demonstrating that the two frameworks are not
alternative descriptions but complementary projections of the
same underlying geometry.
the conditioned legacy figure gallery in Appendix~\ref{app:conditioned-legacy-figure-gallery} summarises the coverage
progression from MES algebraic bounds through the defect-variable
framework of this work to a full Bayesian template analysis.


\subsection{The DCP metric and parameter identification}

The Dipole Cosmological Principle (DCP) of Krishnan, Mondol \&
Sheikh-Jabbari~\cite{Krishnan2022,Krishnan2023} is the locally
rotationally symmetric, tilted Bianchi~V cosmology with metric
\begin{equation}\label{eq:DCP-metric}
ds^2 = -dt^2 + X^2(t)\,dz^2
     + Y^2(t)\,e^{-2A_0 z}(dx^2 + dy^2)\,,
\end{equation}
where $X(t)$ and $Y(t)$ are the longitudinal and transverse
scale factors and $A_0 > 0$ encodes spatial curvature.
Decomposing $X = ae^{2b}$, $Y = ae^{-b}$ with volumetric
scale factor $a = (XY^2)^{1/3}$ and mean Hubble parameter
$H = \dot{a}/a$, the translation proceeds through three exact
identifications (Appendix~\ref{app:DCP}):
\begin{align}
\Sigstd &= \frac{\dot{b}^2}{H^2}
= \frac{\sigma_{\mathrm{DCP}}^2}{9H^2}
=: \Omega_\sigma^{(\mathrm{DCP})}\,,
\label{eq:Sig-DCP}\\
\Otilt &= (1+w)\,\Omega_m\,\sinh^2\!\beta\,,
\label{eq:Otilt-DCP}\\
\Omega_K &= \frac{4A_0^2}{9H^2 X^2}\,.
\label{eq:OmK-DCP}
\end{align}
The defect standardised shear is numerically identical to the DCP
shear density parameter.  The tilt density parameter is the same
expression used in the DCP sum rule.  The curvature is negative
in the convention of Section~\ref{sec:BV-geom}, giving open
spatial sections with the open comparator yielding
$\Okaniso = 0$.


\subsection{The momentum constraint as a dictionary entry}

The DCP momentum constraint
(Eq.~\eqref{eq:BV-momentum-DCP}) provides the critical coupling
\begin{equation}\label{eq:mom-dict}
\Sigstd = \frac{[(1+w)\Omega_m]^2}{4|\Omega_K|}\,\beta^2\,,
\end{equation}
fixing the ratio
$\Sigstd/\Otilt = (1+w)\Omega_m/(4|\Omega_K|)$.  For
$|\Omega_K| = 0.001$, the shear-to-tilt ratio is $\sim 79$;
the defect budget is overwhelmingly shear-dominated, with tilt
contributing only $\sim 1.2\%$.  For $|\Omega_K| = 0.005$ the
ratio drops to $\sim 16$ and the tilt fraction rises to
$\sim 6\%$.  This extreme sensitivity to curvature is a DCP
prediction that the defect framework quantifies.


\subsection{Mutual enrichment}

The defect framework and the DCP address complementary aspects
of the same physics.

The defect framework contributes:
type-by-type specialisation (the master identity applies to all
nine Bianchi types, not just type~V); MES algebraic bounds
that translate CMB multipoles into a shear ceiling
independently of the specific Bianchi dynamics; the nonlinear
$\Teff$ corrections ($\sim 0.8\%$ at S3); and the frame-
attribution bias that identifies the acceleration content of
the CMB dipole.

The DCP contributes:
time evolution through four coupled ODEs
(Appendix~\ref{app:DCP}) governing $\dot{\Sigstd}$,
$\dot{\Otilt}$, $\dot{\Omega}_K$, and $\dot{\beta}$;
attractor analysis establishing that the tilt can grow even
while the shear decays under the cosmic no-hair theorem; and
multi-fluid structure (the radiation--matter relative tilt
$\beta_m - \beta_r$ generates additional defect components
invisible to the single-fluid DCP).


\subsection{Five predictions from the combination}

Five predictions emerge only from the combination of both
frameworks, unavailable to either alone.

\emph{Shear-to-tilt diagnostic.}
The ratio
$f_\sigma/f_{\mathrm{tilt}} = (1+w)\Omega_m/(4|\Omega_K|)$
is a geometric observable: an independent measurement of the
shear (via the quadrupole) and tilt (via the dipole)
contributions would test the Bianchi type.

\emph{MES ceiling on DCP initial conditions.}
The MES bounds constrain DCP initial conditions at any epoch
where CMB data are available, providing an algebraic ceiling
that the DCP evolution equations cannot derive independently.

\emph{Redshift-dependent filling fraction.}
Under the DCP tilt instability, $\FF(z)$ exhibits a
plateau-then-upturn profile: approximately constant for
$0.5 \lesssim z \lesssim z_*$ and rising toward $z = 0$---a
falsifiable prediction testable with redshift-binned dipole
measurements from Euclid and DESI.

\emph{Five-order-of-magnitude $\beta$-gap.}
The maximum curvature-sourced tilt is
$\beta_{\max}^{(\mathrm{curv})} \sim 10^{-7}$ (from the DCP
constraint at $|\Omega_K| = 0.002$), while the anomaly
amplitude is $\beta \sim 10^{-3}$---a gap of five orders of
magnitude that the DCP tilt instability must bridge.

\emph{Defect budget reallocation.}
The cosmic no-hair theorem forces $\Sigstd \to 0$
exponentially, but if the DCP tilt instability simultaneously
drives $\Otilt$ upward, the total defect undergoes an internal
transfer from shear to tilt, so that the late-time defect is
overwhelmingly tilt-dominated ($f_{\mathrm{tilt}} \to 1$).


\bigskip
\noindent\emph{Scope note on the Tsagas bridge
(Sections~\ref{sec:disc-tsagas} ff.).}\quad
The material that follows connects the homogeneous defect
framework to the perturbative tilted-FLRW programme of Tsagas
and collaborators.  This connection is a \emph{discussion-level
extrapolation} that lies outside the formal scope of the
algebraic identity.  The homogeneous framework provides the tilt
rapidity $\beta$ and the departure $x$; the Tsagas
formulae---apparent deceleration
$\hat{q}_0 = q_{0,\mathrm{true}} + \Delta q(\beta, d)$,
peculiar Jeans scale $\lambda_J^{\mathrm{pec}}$, directional
Hubble correction $\Delta H/H$---operate in the perturbative
tilted-observer regime at $d \ll \lambda_H$ and require
additional assumptions (survey selection function, calibrator
angular distribution, single-stream flow validity) that have not
been validated against observational data.  The pushforward
posterior $\hat{q}_0 = 207$ with $95\%$ HPD $[154, 260]$ at
$d = 40$ Mpc ($\hat{q}_0 \approx 1.15$ at $d = 200$ Mpc;
the conditioned legacy figure gallery in Appendix~\ref{app:conditioned-legacy-figure-gallery}) is exact given
$\beta$, but the physical interpretation---that this apparent
deceleration ``explains'' measured local-versus-global tension in
$q_0$---requires survey-level verification.  We present these
extrapolations as testable hypotheses, not established
conclusions.\bigskip


% === INLINED: ch08_solver_and_mechanisms.tex ===
% ============================================================
% ch08_solver_and_mechanisms.tex
% New discussion sections for ch08
% Part of: Tetrad-Based Departure Decomposition (Jiwon Park)
% Session S6 draft — 2026-03-30
% ============================================================


\section{Species-resolved legacy transfer validation context}
\label{sec:disc-solver}

The legacy transfer layer (Chapter~\ref{ch:results}) provides a
species-resolved decomposition of the CMB quadrupole power $D_2$
within the departure parameter framework.  We retain this material as
external/proxy transfer provenance and validation context, not as
native low-ell solver validation.  We now discuss the principal results
and their implications for the departure parameter pipeline.


\subsection{Neutrino dominance of the quadrupole}
\label{sec:disc-nu-dom}

The most consequential result of the solver is the
discovery that neutrinos contribute $98.4\%$ of the total
$D_2$ at recombination, with a neutrino-to-photon
quadrupole ratio $N_2/F_2 = 28$
(Section~\ref{sec:nu-dominance}).  The physical mechanism
is the asymmetry between Thomson-suppressed photons and
free-streaming neutrinos under a common shear source
(Proposition~\ref{prop:nu-dominance}): the shear
$\sigma_{ab}$ drives both species' quadrupoles, but
Thomson scattering suppresses the photon response by a
factor $\dot{\tau}/H \sim 40$ before recombination, while
neutrinos accumulate quadrupole power without collisional
damping.

This result has three implications for the departure
framework.  First, the calibrated transfer function $T_2$
that maps $\Sigstd$ to $D_2$ is dominated by the neutrino
channel: $T_2 \approx T_{2,\nu}$ to $98\%$ accuracy,
and the photon contribution is a $1.6\%$ correction.
The VER06 transfer-conditional legacy value
$\ln \mathcal{B}_{\rm tilt} = +25.29 \pm 0.07$ therefore
tracks primarily the neutrino anisotropic stress proxy
$\pi_\nu^{ab}$, not the photon anisotropy that is directly
observed.  Second, the
$\Teff$ closure for neutrinos is exact, not approximate:
since neutrinos are collisionless after decoupling, the
tangency diagnostic $D_{\nu,\geq 2} = 0$
(Eq.~\ref{eq:D-diagnostic}), and the finite-dimensional
$\Teff$ manifold is an invariant submanifold of the full
Boltzmann dynamics.  Third, the neutrino dominance is
robust against the choice of $\ell_{\rm max}$, the
recombination model, and the reionisation optical depth:
the ratio $N_2/F_2$ varies by less than $3\%$ across the
parameter ranges tested by the legacy transfer path.


\subsection{Oracle deviation budget}
\label{sec:disc-oracle}
\label{sec:oracle-deviation}

The FLRW oracle (Section~\ref{sec:flrw-oracle}) provides
an independent validation target.  The oracle computes
$D_2 = 931~\mu\mathrm{K}^2$ for the fiducial Planck~2018
cosmology, which is $8\%$ below the CAMB reference value
of $1018~\mu\mathrm{K}^2$.  The deviation is accounted for
by two identified sources: the analytic transfer function
$\Psi(k, \eta)$ misses the late integrated Sachs--Wolfe
contribution ($\sim\!5\%$), and the absence of neutrino
anisotropic stress feedback on the gravitational potential
accounts for the remaining $\sim\!3\%$.

The oracle is a toy model, not a production tool.  Its
role is to provide a closed-form FLRW-limit baseline for
legacy transfer bookkeeping.  The AniCLASS
code~\cite{VicentePereiraPitrou2026}, compiled from source
and compared against $15$ historical grid points across
Bianchi types VII$_h$ and IX, is retained here as an
external/proxy comparison record.  It is not native low-ell
solver validation.  At the tested grid points, the legacy
transfer path and AniCLASS agree on the type-dependent
transfer $D_2(\Sigma^2)$ to within $5\%$ for BVIIh and to
within the numerical noise floor for BIX.  The AniCLASS
comparison also records the BVIIh cascade effect:
$D_2(\mathrm{BVIIh}) < D_2(\mathrm{BI})$ at fixed
$\Sigma^2$, because the curvature coupling leaks power
from $\ell = 2$ to higher multipoles
(Section~\ref{sec:bviih-cascade}).


\subsection{Linearity and the FLRW recovery}
\label{sec:disc-linearity}

The solver validates the linearity $D_2 \propto \Sigma^2$
over four decades in $\Sigma^2$ ($10^{-12}$ to $10^{-6}$)
with a relative spread below $0.1\%$
(Section~\ref{sec:solver-linearity}).  The FLRW limit is
recovered exactly: at $\Sigma^2 = 0$, all five
observables ($D_2$, $E_2$, $F_\ell$, $N_\ell$, $E_\ell$)
vanish to machine precision
(Section~\ref{sec:flrw-recovery}).  These two properties
together confirm that the solver introduces no spurious
anisotropy and that the transfer function $T_2 = D_2 /
\Sigma^2$ is a well-defined constant in the linear regime.

The reionisation correction $\exp(-2\tau_{\rm reio})$
with $\tau_{\rm reio} = 0.050$ applies uniformly to the
$\ell = 2$ power, and the Bayes factor $\ln \mathcal{B}$
is exactly invariant under reionisation
(Section~\ref{sec:reionisation-invariance}).  The
polarisation feedback---the back-reaction of the E-mode
quadrupole on the temperature quadrupole through the
Thomson polarisation source $\Pi_{ab}$---is $0.006\%$, a
correction that is negligible at any foreseeable precision.


% ============================================================
\section{Mechanism attribution after Clarkson--Maartens}
\label{sec:disc-mechanisms}
% ============================================================

The departure parameter pipeline gives a
legacy transfer-conditional statistical preference for nonzero tilt:
the Bayes
factor $\ln \mathcal{B}_{\rm tilt} = +25.29 \pm 0.07$
measures the statistical weight of the data against the
untilted null hypothesis, without assuming a specific
physical mechanism for the tilt.  We now assess the
theoretical status of the candidate mechanisms as of
March~2026, in light of two developments that have
reshaped the field.


\subsection{The Tsagas fast-growth mechanism: killed}
\label{sec:disc-tsagas-killed}

Clarkson \& Maartens~\cite{ClarksonMaartens2026}
demonstrated that the enhanced peculiar-velocity growth
claimed by Tsagas and
collaborators~\cite{Tsagas2010,TsagasKadiltzoglou2013}---$v
\propto a^{3/2}$ to $a^2$ on super-Hubble scales during
matter domination, compared to the standard $v \propto
a^{1/2}$---arises from treating constraint-determined
quantities as independent sources in the 1+3 covariant
perturbation equations.  When the constraint equations are
imposed consistently, the standard perturbative result
$v \propto t^{1/3}$ (equivalently $v \propto a^{1/2}$)
is recovered exactly.

The Clarkson--Maartens argument is technically decisive: it
identifies the specific algebraic step at which the Tsagas
derivation violates the Einstein constraints, and shows
that restoring consistency eliminates the enhanced growth.
No rebuttal has been filed as of March~2026, despite a
vigorous defence of the fast-growth mechanism by
Past{\'e}n \& Tsagas~\cite{PastenTsagas2026} published
six weeks before the Clarkson--Maartens preprint.

For the departure framework, the consequence is
straightforward: the Tsagas mechanism cannot be invoked to
explain the observed dipole excess through enhanced
peculiar velocities in a perturbed FLRW setting.  The
tilted-observer programme of
Section~\ref{sec:disc-tsagas}---specifically the apparent
deceleration $\tilde{q}$ and its connection to the
$H_0$ tension---remains valid as a kinematic framework,
but the dynamical mechanism (fast peculiar-velocity
growth) that was supposed to drive $\tilde{q}$ away
from $q$ is no longer available.


\subsection{The Krishnan--Mondol tilt instability: survives}
\label{sec:disc-krishnan}

The dipole cosmology programme of Krishnan, Mondol \&
Sheikh-Jabbari~\cite{Krishnan2023} operates in a
fundamentally different mathematical setting from the
Tsagas programme.  Where Tsagas works within perturbed
FLRW using the 1+3 covariant formalism, Krishnan works
with exact Bianchi~V solutions described by a
finite-dimensional ODE system.  The Clarkson--Maartens
critique targets inconsistencies in the perturbative
truncation, not in the exact dynamics; the Krishnan
framework is therefore unaffected.

The key result of the dipole cosmology programme is that
the cosmic no-hair theorem forces metric shear
$\sigma \to 0$ during $\Lambda$-dominated expansion, but
does \emph{not} force the fluid tilt $\beta \to 0$.  The
tilt rapidity can grow or persist at late times, driven by
the shear--tilt decoupling that occurs when
$\Lambda$ dominates the expansion.  In the multi-fluid
extension~\cite{Krishnan2023}, the relative flow between
radiation and matter generically grows at late times---a
behaviour that naturally explains why the matter dipole
(traced by galaxy surveys) exceeds the CMB dipole.

For the departure framework, the Krishnan tilt instability
provides a viable dynamical mechanism that is independent
of the Tsagas programme and compatible with the
phenomenological detection.  The tilt instability does not
by itself predict the amplitude $\beta \sim 10^{-3}$:
that requires specifying initial conditions at decoupling,
which is a free parameter in the exact Bianchi~V system.
The departure pipeline constrains $\beta$ from the data;
the Krishnan framework provides the theoretical context
in which a nonzero $\beta$ is dynamically natural.


\subsection{The khronon field as sole surviving standard source}
\label{sec:disc-khronon}

Mart{\'i}n, Skordis, Bartlett, Desmond, Ferreira \&
Yasin~\cite{MartinSkordis2025} systematically evaluated
four candidate tilt sources within Bianchi~V: spatial
curvature, heat flow, electromagnetic fields, and a
khronon field (the hypersurface-orthogonal limit of
Einstein-aether theory).  Three of the four fail to
produce the observed dipole amplitude by five to thirteen
orders of magnitude under current observational constraints.
The curvature-driven mechanism---the one most closely
related to the departure parameter
$\Okaniso$---yields $\beta \sim 10^{-13}$ for dust,
far below the observed $\beta \sim 10^{-3}$.

Only the khronon field, acting as dust-like dark matter
with an independent four-velocity, achieves
$\beta \sim 2$--$4 \times 10^{-3}$ within the
post-GW170817 parameter space.  The viable parameter
window requires $\Omega_K \sim \Omega_f$ (the khronon
constitutes a significant fraction of the total matter
density) and adiabatic initial conditions with
$\delta_K \sim 5 \times 10^{-6}$, set by the CMB
quadrupole shear bound.

This result does not invalidate the departure parameter
detection; it constrains the interpretation.  The Bayes
factor $\ln \mathcal{B}_{\rm tilt} = +25.29$ measures the
phenomenological signal for nonzero tilt across seven
data channels.  The Mart{\'i}n--Skordis analysis narrows
the physical mechanism that could produce this signal.
Detection and mechanism attribution are separable
questions: the pipeline establishes the former; the
theoretical analysis constrains the latter.

\begin{remark}[Detection versus mechanism]
\label{rem:detection-mechanism}
The departure parameter framework is mechanism-agnostic
by construction: the master identity
(Eq.~\ref{eq:master-identity}) is an algebraic
relation among geometry-frame kinematic quantities, and
the Bayesian evidence computation does not assume a
specific tilt source.  The Mart{\'i}n--Skordis result
that only the khronon field achieves $\beta \sim
10^{-3}$ under standard ancillary constraints narrows
the interpretation of the detection but does not
weaken the detection itself.
\end{remark}


\subsection{The separation principle}
\label{sec:disc-separation}

The three developments---Clarkson--Maartens killing Tsagas,
Krishnan surviving independently, and Mart{\'i}n--Skordis
narrowing the source to the khronon---establish a clean
separation between three levels of the analysis:
\emph{(i)}~the legacy transfer-conditional statistical preference
for nonzero tilt
(this work, $\ln \mathcal{B} = +25.29$);
\emph{(ii)}~the dynamical viability of persistent tilt at
late times (Krishnan tilt instability, exact Bianchi~V);
and \emph{(iii)}~the microscopic source that generates the
tilt (khronon field, or competing exotic mechanisms such
as QCD axion-induced isocurvature).

The departure framework operates at level~(i) and provides
the observational anchor for levels~(ii) and~(iii).  The
species-resolved solver of
Section~\ref{sec:multi-species-teff} adds a diagnostic
capability: by tracking each species' quadrupole
independently, the solver can distinguish tilt signatures
that affect all species equally (as in the curvature-driven
mechanism) from those that create differential tilt between
radiation and matter (as in the Krishnan multi-fluid
instability or the khronon field).  This distinction is not
accessible to the scalar-amplitude pipeline of VER06 but
is a primary target of the direction-dependent programme
(Section~\ref{sec:solver-programme}).


% ============================================================
\section{Multi-fluid implications for the framework}
\label{sec:disc-multi-fluid}
% ============================================================

The species-resolved framework of
Section~\ref{sec:multi-fluid} extends the single-fluid
departure decomposition to a realistic cosmological
medium.  We now discuss the implications of this extension
for the methodological limitations identified in
Section~\ref{sec:disc-limitations}.


\subsection{From single-fluid to species-resolved tilt}
\label{sec:disc-species-tilt}

The safe-route tilt bound
$\beta \leq \varepsilon_1/(1 + \eta_{\dot{u}})$
(Section~\ref{sec:corrected-bound}) was derived under the
single-barotropic-fluid assumption, which gives
$\eta_{\dot{u}} = 1/12$ through the tilt--acceleration
relation $\dot{u}_a = [w/(3(1+w))]\,\tilde{D}_a \ln
\rho$.  In a multi-component universe, the
four-acceleration $\dot{u}_a$ retains independent degrees
of freedom not captured by the tilt rapidity alone
(Section~\ref{sec:disc-limitations}).

The species-resolved framework
(Section~\ref{sec:multi-fluid}) provides the natural
language for this extension.  Each species carries its own
tilt through the dipole $\Theta^{(s)}_a$
(Eq.~\ref{eq:dipole-drift}), and the total
four-acceleration receives contributions from all species
through the Einstein equations.  The single-fluid
$\eta_{\dot{u}} = 1/12$ is replaced by a species-weighted
effective correction that depends on the relative tilts
of radiation, matter, and dark energy.  Deriving the
explicit multi-fluid bound requires specifying the
inter-species tilt angles---additional free parameters that
are not constrained by current data---and is deferred to
the direction-dependent programme.

For the present scalar-amplitude analysis, the single-fluid
bound remains the tightest (most constraining) available,
and the frame-attribution bias $\eta_{\dot{u}} \approx
8.3\%$ reported in Section~\ref{sec:frame-bias} is a
conservative upper bound on the multi-fluid correction.


\subsection{The local-patch connection}
\label{sec:disc-local-patch}

The peculiar velocity $v_{\rm loc}$ of the Local Group
relative to the CMB rest frame enters the departure
analysis through the observer-frame boost
(Section~\ref{sec:frame-problem}).  In the species-resolved
framework, this local motion is encoded in the photon
dipole $\Theta^{(\gamma)}_a$ and affects only the
temperature power spectrum (TT) at $\ell = 1$; it does not
source E-mode polarisation.  Cosmological tilt, by
contrast, affects both TT and EE through the pre-recombination
anisotropic stress
(Sections~\ref{sec:thomson-coupling}--\ref{sec:nu-dominance}).

This TT-vs-EE distinction provides a clean separation
between local and cosmological contributions to the
observed dipole.  The Bayesian filling fraction
$\mathcal{F}_{\rm Bayes} = 0.093 \pm 0.025$
(Section~\ref{sec:filling-fraction}) can be decomposed
into a cosmological component (constrained by the
multi-channel evidence) and a local component (absorbed by
the observer-frame boost).  The CosmicFlows-4 bulk flow
of $419 \pm 36$~km/s at
$200\,h^{-1}$~Mpc~\cite{Watkins2023} is consistent with
the local component but exceeds the $\Lambda$CDM
expectation by $\sim\!3\sigma$, suggesting that part of
the signal may have a cosmological origin.

The species-resolved solver can quantify this
decomposition by computing the differential response of
$D_2$ to local versus cosmological tilt sources.  In the
linear regime, the two contributions add in quadrature:
\begin{equation}\label{eq:D2-local-cosmo}
D_2 = D_{2,\rm cosmo}(\Sigstd, \beta)
    + D_{2,\rm local}(v_{\rm loc}, W_R)
    + O(\text{cross-terms})\,,
\end{equation}
where $W_R$ is the local-patch window function
(Section~\ref{sec:quadratic-sources}).  The cross-terms
vanish at linear order because the cosmological shear and
the local velocity source different angular
patterns on the sky.  At second order, the cross-terms
(Eq.~\ref{eq:local-patch-2nd}) are of order
$W_R\,(\Sigma^2)^{1/2}\,(v_{\rm loc}/c)
\sim 10^{-6}$, which is negligible.


\subsection{Implications for the $H_0$ tension}
\label{sec:disc-H0}

The tilted Hubble rate
(Eq.~\ref{eq:tilted-H-tsagas}) predicts a systematic
offset between the Hubble constant measured by a tilted
observer and the background expansion rate:
$H_{\rm tilt} / H_{\rm FLRW} = 1 + \tanh^2\!\beta\,
\Sigstd + O(\beta^4)$.  For the VER06 legacy numerical values
($\beta = 1.360 \times 10^{-3}$,
$\Sigstd \sim 10^{-5}$), the fractional offset is
\begin{equation}\label{eq:H0-offset}
\frac{\Delta H_0}{H_0}
\sim \beta^2\,\Sigstd
\sim (1.4 \times 10^{-3})^2 \times 10^{-5}
\sim 2 \times 10^{-11}\,,
\end{equation}
which is many orders of magnitude below the $\sim\!8\%$
$H_0$ discrepancy between the Planck CMB value
($67.4 \pm 0.5$~km/s/Mpc) and the SH0ES distance-ladder
value ($73.0 \pm 1.0$~km/s/Mpc).

The departure framework therefore does not resolve the
$H_0$ tension through the tilted Hubble rate mechanism
alone---the detected tilt is too small by nine orders of
magnitude.  This negative result is itself informative: it
rules out the possibility that the observed cosmic dipole
anomaly and the $H_0$ tension share a common origin in a
global cosmological tilt, at least within the Bianchi
framework and at the currently detected amplitude.

Alternative routes through which anisotropy could affect
$H_0$ measurements---such as direction-dependent Hubble
diagrams, anisotropic void models, or backreaction-induced
scale dependence (Section~\ref{sec:buchert-equations})---remain
open and are structurally distinct from the global tilt
mechanism.  The species-resolved solver provides the
infrastructure to test these alternatives once
direction-dependent data become available.

% === END INLINED ===

% === INLINED: ch08_sec83_tsagas_expanded.tex ===
% ============================================================
% ch08_sec83_tsagas_expanded.tex
% Expanded §8.3: The Tsagas tilted-observer programme
% Replaces existing sec:disc-tsagas (lines 299–366 of ch08)
% Depends on: ch02_sec7 (§2.7), TF-T03/T06 results, TF-R03/R04 data
% ============================================================

\section{The Tsagas tilted-observer programme}
\label{sec:disc-tsagas}

Our goal in this section is to connect the homogeneous defect
framework developed in Chapters~\ref{ch:framework}--\ref{ch:results}
to the perturbative tilted-FLRW programme of
Tsagas~\cite{Tsagas2010,TsagasKadiltzoglou2013,Tsagas2021_Jeans},
which provides the dynamical mechanism that the algebraic defect
identity lacks: a relativistic explanation for why observers
embedded in bulk peculiar flows perceive different kinematics from
the background Hubble flow.  Section~\ref{sec:tilted-kinematics}
established the formal bridge; here we confront the theoretical
predictions with observational data across five independent probes.


% ────────────────────────────────────────────────────────────
\subsection{The tilted-FLRW programme (2010--2026)}
\label{sec:disc-tsagas-programme}

The programme builds on a single physical insight: the energy flux
$q_a = (\rho + p)v_a$ carried by a tilted observer gravitates in
general relativity but has no Newtonian counterpart.  This term
enters the tilted Raychaudhuri equation through
$D^a\hat{A}_a$---the divergence of the matter-frame
four-acceleration---and modifies the deceleration parameter
measured by any observer not at rest with respect to the Hubble flow.
The master relation (Section~\ref{sec:tilted-raych-defect},
equation~\eqref{eq:delta-q-exact}) is
\begin{equation}\label{eq:disc-qtilde}
\tilde{q}
= q - \frac{1}{3H^2}\,D^a\hat{A}_a\,.
\end{equation}
On sub-Hubble scales ($d \ll \lambda_H$), the self-consistent
closure $|\vartheta|/H = \beta\lambda_H/d$
(equation~\eqref{eq:vartheta-estimate}; conditional on a single
coherent monopolar peculiar flow) gives the defect-variable
form
\begin{equation}\label{eq:disc-Dq}
\Delta q(\beta, d)
= \frac{\beta}{9}
  \Bigl(\frac{\lambda_H}{d}\Bigr)^{\!3}\,,
\end{equation}
where $\beta$ is the tilt rapidity constrained by the defect
identity and $d$ is the comoving survey depth.  The cubic scale
dependence produces enormous amplification: at
$d = 100\;\mathrm{Mpc}$, the factor
$(\lambda_H/d)^3 \approx 8.8 \times 10^4$; at
$d = 300\;\mathrm{Mpc}$, it drops to
$\sim 3.3 \times 10^3$.

\emph{Velocity growth.}
The same energy-flux feedback drives the peculiar velocity
growth law.  Tsagas \& Kadiltzoglou~\cite{TsagasKadiltzoglou2013}
showed that, on an Einstein--de~Sitter background, the growth
index satisfies $p^2 - (1/3)p - 2/3 = 0$, giving $p = 1$ and
$v \propto a^{3/2}$ (the ``GR minimum''); the
long-wavelength limit gives $v \propto a^2$
(Section~\ref{sec:tilted-raych-defect},
Remark~\ref{rem:scale-hierarchy}).
Both exceed the Newtonian $v \propto a^{0.69}$
($p = 0.457$ from $p^2 + p - 2/3 = 0$) by a growth-index
ratio of $\sim\!2$--$4$.  The cumulative enhancement from
matter--radiation equality to $z = 0$ is a factor of
$\sim\!750$, qualitatively explaining why the CF4 bulk flow at
$200\,h^{-1}\;\mathrm{Mpc}$ exceeds $\Lambda$CDM predictions
at $4.8\sigma$~\cite{Watkins2023}.

\emph{The peculiar Jeans length.}
The scale at which the apparent deceleration changes sign defines
$\lambda_J^{\mathrm{pec}}
= \lambda_H(\beta/(9q))^{1/3}$
(equation~\eqref{eq:lambda-J-defect}; conditional on the
self-consistent closure and a fixed background $q$).  At the S3 anomaly
amplitude ($\beta = 1.36 \times 10^{-3}$, $q = 1/2$):
$\lambda_J^{\mathrm{pec}} \approx 298\;\mathrm{Mpc}$
(Table~\ref{tab:disc-tilted-H0}), coinciding with the CF4 measurement
depth.  Below this scale, tilted observers in contracting flows
perceive $\tilde{q} < 0$; above it, all observers agree on $q$.


% ────────────────────────────────────────────────────────────
\subsection{Apparent deceleration as a defect observable}
\label{sec:disc-apparent-q}

The directional structure of~\eqref{eq:disc-Dq} is dipolar:
$\tilde{q}(\hat{n}) = q + \Delta q_{\mathrm{dip}}\cos\theta$,
where $\theta$ is the angle between the line of sight and the
tilt direction.  Three falsifiable predictions follow:
\emph{(a)}~the dipole is aligned with the CMB dipole direction;
\emph{(b)}~the amplitude decays with redshift; and
\emph{(c)}~the monopole $q_m$ is consistent with zero
acceleration.

\emph{Colin et~al.\ (2019).}
A maximum-likelihood fit to 740 JLA SNe Ia parametrised as
$q_0(\hat{n}) = q_m + q_d\,(\hat{n}\cdot\hat{d})\,e^{-z/S}$
yields $q_d = -8.03$, $S = 0.0262$, $q_m = -0.157$ at
$3.9\sigma$ rejection of isotropy~\cite{Colin2019}.  All three
Tsagas predictions are confirmed: dipole aligned with the CMB
direction, exponential redshift decay (characteristic depth
$d_S = cS/H_0 = 117\;\mathrm{Mpc}$), and $q_m$ consistent with
zero at $1.4\sigma$.

\emph{Translation to $\beta$.}
Equating the Colin amplitude at a reference redshift $z_{\mathrm{ref}}$
to the Tsagas formula~\eqref{eq:disc-Dq}, the tilt rapidity is
\begin{equation}\label{eq:disc-betaSNe}
\beta_{\mathrm{SNe}}
= 9\,|q_d|\,z_{\mathrm{ref}}^3\,e^{-z_{\mathrm{ref}}/S}\,.
\end{equation}
At $z_{\mathrm{ref}} = 0.05$ ($d = 222\;\mathrm{Mpc}$, the JLA
median depth):
$\beta_{\mathrm{SNe}} = 1.34 \times 10^{-3}$, consistent
with the CF4 value
$\beta_{\mathrm{CF4}} = (1.334 \pm 0.267) \times 10^{-3}$
within the $\sim\!20\%$ CF4 uncertainty
(Table~\ref{tab:colin-beta}; the conditioned legacy figure gallery in Appendix~\ref{app:conditioned-legacy-figure-gallery}).  The agreement is
$z_{\mathrm{ref}}$-dependent: at $z = 0.04$ the ratio
$\beta_{\mathrm{SNe}}/\beta_{\mathrm{CF4}} = 0.75$; at $z = 0.06$
it is $1.18$.  Both the $\beta$ translation and the Jeans length
evaluation share the self-consistent closure
(equation~\eqref{eq:vartheta-estimate}); the match therefore
tests the Tsagas $d^{-3}$ scaling against the Colin exponential
parametrisation, not the absolute tilt amplitude.

\begin{table}[t]
\centering
\footnotesize
\setlength{\tabcolsep}{4pt}
\caption{Translation of the Colin et~al.\ (2019) dipolar
deceleration to the tilt rapidity $\beta$ via the Tsagas
formula~\eqref{eq:disc-Dq}, evaluated at six reference
redshifts.  The CF4 value
$\beta_{\mathrm{CF4}} = 1.334 \times 10^{-3}$ is shown
for comparison.}
\label{tab:colin-beta}
\begin{tabular}{ccccl}
\toprule
$z_{\mathrm{ref}}$ & $d_{\mathrm{ref}}$ (Mpc) &
  $|\Delta q|_{\mathrm{Colin}}$ &
  $\beta_{\mathrm{SNe}}$ &
  $\beta_{\mathrm{SNe}}/\beta_{\mathrm{CF4}}$ \\
\midrule
0.02 &  89 & 3.74 & $2.7 \times 10^{-4}$ & 0.20 \\
0.03 & 133 & 2.56 & $6.2 \times 10^{-4}$ & 0.47 \\
\textbf{0.05} & \textbf{222} & \textbf{1.19} &
  $\mathbf{1.34 \times 10^{-3}}$ & \textbf{1.00} \\
0.08 & 356 & 0.38 & $1.75 \times 10^{-3}$ & 1.31 \\
0.10 & 445 & 0.18 & $1.59 \times 10^{-3}$ & 1.19 \\
0.15 & 667 & 0.03 & $7.9 \times 10^{-4}$  & 0.60 \\
\bottomrule
\end{tabular}
\end{table}

\emph{Sah et~al.\ (2025).}
The analysis of 1533 Pantheon+ SNe
Ia reports a dipolar $H_0$ variation exceeding
$1.5\;\mathrm{km\,s}^{-1}\mathrm{Mpc}^{-1}$ in the shell
$0.023 < z < 0.15$ at $>5\sigma$ in all reference frames
(heliocentric, CMB, Local Group)~\cite{Sah2025}.  The Tsagas
formula predicts $\Delta H \approx 1.3\;\mathrm{km\,s}^{-1}
\mathrm{Mpc}^{-1}$ at the inner edge of the Sah shell
($z = 0.023$, $d \approx 102\;\mathrm{Mpc}$)---marginally
below the reported threshold but order-of-magnitude consistent.

\emph{Channel independence.}
Does the dipolar deceleration provide a new evidence channel?
The corrected Fisher information comparison (TF-N02) gives
$I_i/I_c \approx 0.6$ at $z = 0.05$---\emph{comparable} to the
CF4 channel, not subdominant.  The decisive factor is the high
physical correlation: both channels probe the same $\beta$ at
similar depths, giving $\rho_{ci} \gtrsim 0.8$ and
$\Delta\ln\mathcal{B} < 0.1$.  We therefore do not add
channel~(i) to the evidence computation.  The dipolar deceleration
is a powerful \emph{consistency check} and a \emph{test of
scale dependence} ($\Delta q \propto d^{-3}$), not an
independent parameter constraint.


% ────────────────────────────────────────────────────────────
\subsection{The progenitor age-bias route}
\label{sec:disc-age-bias}

An independent astrophysical programme arrives at the same
qualitative conclusion---the present universe may not be
accelerating---through an entirely different mechanism.

\emph{The Yonsei result.}
Chung et~al.~\cite{Chung2025} establish a $5.5\sigma$ correlation
between post-standardisation Hubble residual and host-galaxy
stellar population age at $s = -0.030 \pm 0.004\;\mathrm{mag\,
Gyr}^{-1}$ from $>300$ hosts across $z < 0.45$.
Son et~al.~\cite{Son2025} apply this correction to the
Pantheon+ and DES-SN5YR catalogues, combined with DESI~DR2 BAO
and Planck+ACT CMB data.  DESI BAO + CMB alone
(no supernovae) prefer a $w_0w_a$CDM cosmology with
$w_0 = -0.42$, $w_a = -1.75$, $\Omega_m = 0.353$, yielding
$q_0 = +0.092 \pm 0.20$---marginally positive (decelerating)
and inconsistent with $\Lambda$CDM at $3.1\sigma$.  After
age-bias correction, the combined SN+BAO+CMB analysis shows
$> 9\sigma$ tension with $\Lambda$CDM.

We note that the Son et~al.\ $q_0$ is obtained within the
$w_0w_a$CDM parametrisation; alternative dark-energy
parametrisations or the standard $\Lambda$CDM model return
$q_0 = -0.55$.  The shift in $q_0$ is therefore
model-dependent at the $\sim\!0.6$ level.

\emph{The three-term decomposition}
(the conditioned legacy figure gallery in Appendix~\ref{app:conditioned-legacy-figure-gallery}).
In the most general case, the observed deceleration parameter
from SNe Ia contains three contributions:
\begin{equation}\label{eq:disc-q-decomp}
q_0^{(\mathrm{obs})}
= q_0^{(\mathrm{true})}
  + \Delta q^{(\mathrm{tilt})}(\beta, d, \hat{n})
  + \Delta q^{(\mathrm{age})}(z)\,.
\end{equation}
Standard $\Lambda$CDM sets both corrections to zero, yielding
$q_0 = -0.55$.  The Tsagas programme sets
$\Delta q^{(\mathrm{age})} = 0$ and attributes the full offset to
tilt.  The Yonsei programme sets
$\Delta q^{(\mathrm{tilt})} = 0$ and attributes it to age bias.
If both operate, the budget must close:
$\Delta q^{(\mathrm{tilt})} + \Delta q^{(\mathrm{age})}
\approx 0.64$ for the BAO+CMB prediction to hold.

At $d = 222\;\mathrm{Mpc}$, the Tsagas formula gives
$\Delta q^{(\mathrm{tilt})} \approx 1.19$, exceeding the total
required offset by a factor $\sim\!1.9$.  This overproduction
reflects the breakdown of the self-consistent closure at
$d \lesssim 300\;\mathrm{Mpc}$ :
the $d^{-3}$ formula overpredicts the Colin amplitude by
$\sim\!2\times$ at these depths.  At $d \sim 350$--$400\;
\mathrm{Mpc}$, $\Delta q^{(\mathrm{tilt})} \sim 0.3$--$0.5$,
and the budget closes with a residual
$\Delta q^{(\mathrm{age})} \sim 0.15$--$0.35$.

\emph{Distinguishability.}
The two contributions are partially degenerate but
separable through five observational discriminants
(Table~\ref{tab:disc-discriminants}).

\begin{table}[t]
\centering
\footnotesize
\setlength{\tabcolsep}{3pt}
\caption{Observational discriminants between the age-bias
(Route~A) and tilt (Route~B) contributions to the apparent
acceleration.  Each test exploits a qualitative difference in the
predicted $q_0$ residual pattern.}
\label{tab:disc-discriminants}
\begin{tabular}{lcc}
\toprule
Test & Route A (age bias) & Route B (tilt) \\
\midrule
Directional dependence & Isotropic & Dipolar \\
Redshift scaling & Grows with $z$ & Falls as $d^{-3}$ \\
Evolution-free subsample & Signal vanishes & Signal persists \\
BAO confirmation & Independent test & Independent test \\
Host-age--$\hat{n}$ cross & Correlated & Uncorrelated \\
\bottomrule
\end{tabular}
\end{table}

The fifth row is the sharpest discriminant.  If the age bias
dominates, the Hubble residual should correlate with host age
but not with $\cos\theta$.  If the tilt dominates, the residual
should correlate with direction but not host age.  A bivariate
fit in (host age, $\cos\theta$) from the Rubin LSST sample
($>20{,}000$ SN host galaxies with measured ages and full-sky
coverage) can break the degeneracy on the timescale of 2028--2030.

\emph{Peculiar Jeans length at $q = 0.09$.}
If the Son et~al.\ $q_0^{(\mathrm{true})} = +0.09$ applies,
$\lambda_J^{\mathrm{pec}} = 525\;\mathrm{Mpc}$---$77\%$ larger
than the EdS value ($298\;\mathrm{Mpc}$) because a smaller
$q$ is easier to overcome.  The contamination zone extends well
beyond the CF4 scale, consistent with the Sah et~al.\ finding
that the SN dipole persists out to $z \sim 0.15$
($d \sim 670\;\mathrm{Mpc}$).


% ────────────────────────────────────────────────────────────
\subsection{The tilted Hubble rate and the $H_0$ tension}
\label{sec:disc-tilted-H0}

The tilted Hubble rate $\hat{H} = H\cosh\beta$ produces a
fractional correction $\Delta H/H = \cosh\beta - 1 \approx
\beta^2/2 \approx 8.9 \times 10^{-7}$ in the homogeneous Bianchi
framework (equation~\eqref{eq:H-ratio-defect}).  The $H_0$ tension
$\Delta H = 5.6\;\mathrm{km\,s}^{-1}\mathrm{Mpc}^{-1}$
($8.4\%$) is five orders of magnitude beyond this reach.

The perturbative Tsagas extension
(equation~\eqref{eq:tilted-H-tsagas}; conditional on the
same self-consistent closure) introduces the sub-Hubble
amplification:
\begin{equation}\label{eq:disc-dH}
\frac{\hat{H} - H}{H}
\approx \frac{\beta}{3}\,\frac{\lambda_H}{d}\,,
\end{equation}
which is scale-dependent and can be significant at local depths.
Table~\ref{tab:disc-tilted-H0} and the conditioned legacy figure gallery in Appendix~\ref{app:conditioned-legacy-figure-gallery} evaluate this correction for
the characteristic depths of the major $H_0$ measurement
programmes.

\begin{table}[t]
\centering
\footnotesize
\setlength{\tabcolsep}{3pt}
\caption{Tilted Hubble correction $\Delta H$ at the effective
depth of various $H_0$ methods, using $\beta = 1.334 \times
10^{-3}$ (CF4).  Values give the \emph{maximum} directional
correction ($\cos\theta = 1$); the sky-averaged contribution
depends on the angular distribution of calibrators.
The SH0ES--Planck tension is
$\Delta H_0 = 5.6\;\mathrm{km\,s}^{-1}\mathrm{Mpc}^{-1}$
($8.4\%$).}
\label{tab:disc-tilted-H0}
\begin{tabular}{lcccc}
\toprule
Scale & $d$ (Mpc) & $\Delta H/H$ &
  $\Delta H$ (km/s/Mpc) &
  \% of tension \\
\midrule
Cepheid hosts    &  30 & 6.6\% & 4.4  & 79\% \\
SH0ES median     &  50 & 4.0\% & 2.7  & 47\% \\
TRGB (Freedman)  &  75 & 2.6\% & 1.8  & 32\% \\
CF4 inner        & 100 & 2.0\% & 1.3  & 24\% \\
CF4 outer        & 200 & 1.0\% & 0.7  & 12\% \\
$\lambda_J^{\mathrm{pec}}$ & 300 & 0.7\% & 0.4 & 8\% \\
BAO low-$z$      & 500 & 0.4\% & 0.3  & 5\% \\
CMB$^{\dagger}$  & 4450 & $9 \times 10^{-5}\%$ &
  $6 \times 10^{-5}$ & $\sim 0$ \\
\bottomrule
\multicolumn{5}{l}{\scriptsize $^{\dagger}$At $d \sim \lambda_H$,
equation~\eqref{eq:disc-dH} is invalid; the Bianchi result
$\beta^2/2$ applies.}
\end{tabular}
\end{table}

At the median SH0ES depth ($d \approx 50\;\mathrm{Mpc}$),
$\Delta H \approx 2.7\;\mathrm{km\,s}^{-1}\mathrm{Mpc}^{-1}$
covers $47\%$ of the tension.  For the TRGB method
(Freedman et~al.\ 2024:
$H_0 = 69.8 \pm 1.3\;\mathrm{km\,s}^{-1}\mathrm{Mpc}^{-1}$),
the tension with Planck is only
$2.4\;\mathrm{km\,s}^{-1}\mathrm{Mpc}^{-1}$, which the tilt
correction at $d \approx 50\;\mathrm{Mpc}$ can in principle
account for (the sky-averaged contribution depends on the angular
distribution of TRGB calibrators, which has not been evaluated).

The directional signature is dipolar:
\begin{equation}\label{eq:disc-H-dir}
\hat{H}(\hat{n})
= H\Bigl[1
  + \frac{\beta}{3}\,\frac{\lambda_H}{d}\,\cos\theta\Bigr]\,.
\end{equation}
At $d = 50\;\mathrm{Mpc}$, the full-sky spread is
$\Delta\hat{H} \approx 5.3\;\mathrm{km\,s}^{-1}
\mathrm{Mpc}^{-1}$---comparable to the tension itself.  The
Migkas et~al.\ (2021) detection of a $9\%$ directional $H_0$
variation in X-ray clusters, with maximum at
$(l, b) \approx (280^\circ, -15^\circ)$ (only $19^\circ$ from
the CF4 axis), is precisely the predicted pattern~\cite{Migkas2021}.

\emph{Conservative, progressive, and data-honest readings.}
The homogeneous Bianchi framework predicts
$\Delta H/H \approx 9 \times 10^{-7}$; the $H_0$ tension cannot
be addressed at this level.  The perturbative formula gives
$2$--$4\%$ at SH0ES depth, but this contribution lies outside the
formal scope of the homogeneous defect identity.  The data-honest
position is that the tilt contribution is scale-dependent: at
CMB scales, negligible ($6 \times 10^{-5}\;\mathrm{km\,s}^{-1}
\mathrm{Mpc}^{-1}$); at SH0ES scales, comparable to the tension.
The exact sky-averaged contribution depends on the angular
distribution of distance-ladder calibrators, which has not been
evaluated for the SH0ES sample.

Three testable predictions distinguish the tilt contribution from
other systematics.
\emph{(a)}~$H_0$ should exhibit a $\cos\theta$ dipole toward
the CF4 direction, decreasing as $d^{-1}$.
\emph{(b)}~The tension should decrease monotonically with
calibrator depth.
\emph{(c)}~The directional $H_0$ maximum should coincide with
the dipolar deceleration maximum.

We frame the perturbative $H_0$ contribution as a
\emph{discussion item}, not a principal result.  The defect
framework's contribution to the $H_0$ question is indirect:
it establishes $\beta$ from CF4, provides the theoretical ceiling
$\betasr \leq \eps_1/(1+\etaud)$ from the CMB dipole, and
derives $\etaud$ as the bridge between the Bianchi and Tsagas
regimes.  The quantitative $H_0$ evaluation requires the
perturbative extension and the SH0ES angular selection
function---both outside the formal scope of this work.


% === INLINED: ch08_H0_selfconsistency.tex ===
% Section content is integrated inline in ch08_discussion.tex

% === END INLINED ===


% ────────────────────────────────────────────────────────────
\subsection{Observational scorecard}
\label{sec:disc-scorecard}

Table~\ref{tab:disc-scorecard} compiles the observational
status of the six independent probes that test Tsagas predictions.

\begin{table}[t]
\centering
\footnotesize
\setlength{\tabcolsep}{3pt}
\caption{Observational scorecard for the Tsagas tilted-FLRW
programme.  ``Direction'' gives galactic coordinates
$(l, b)$ of the detected anisotropy axis.}
\label{tab:disc-scorecard}
\begin{tabular}{lcccl}
\toprule
Probe & Significance & Direction &
  Tsagas prediction & Status \\
\midrule
Sah Pantheon+ dipolar $q_0$
  & $>5\sigma$ & $\sim 7^\circ$ from CMB
  & $\Delta q = (\beta/9)(\lambda_H/d)^3$ & Pattern consistent$^{\ddagger}$ \\
Colin SN dipolar $q$
  & $3.9\sigma$ & $(264^\circ, 48^\circ)^{\dagger}$
  & $\Delta q \propto d^{-3}$ & Pattern consistent$^{\ddagger}$ \\
B\"ohme multi-survey radio
  & $5.4\sigma$ & CMB-aligned
  & Excess dipole & Consistent \\
CF4++ bulk flow
  & $>5\sigma$ & $(282^\circ, +6^\circ)$
  & $v \propto a^{3/2}\text{--}a^2$ & Consistent \\
Secrest quasar dipole
  & $4.9\sigma$ & $(244^\circ, 46^\circ)$
  & Excess dipole & Consistent \\
Migkas X-ray clusters
  & $5.4\sigma$ & $(280^\circ, -15^\circ)$
  & Directional $H_0$ & Pattern consistent$^{\ddagger}$ \\
\bottomrule
\multicolumn{5}{l}{\scriptsize $^{\dagger}$Fixed to CMB dipole
direction in the Colin analysis.}\\
\multicolumn{5}{l}{\scriptsize $^{\ddagger}$Qualitative pattern
matches; specific scaling not yet directly tested against
alternatives.}
\end{tabular}
\end{table}

The strongest new entry is Sah, Rameez, Sarkar \&
Tsagas~\cite{Sah2025}, who report a $>5\sigma$ dipolar
deceleration parameter $|q_d| \approx 8$ at
$S = 0.0262$ ($d \approx 116\;\mathrm{Mpc}$) from the
Pantheon+ compilation (1,533 supernovae).  Their dipole
direction lies within $\sim 7^\circ$ of the CMB dipole
axis, and the monopole $q_m$ is consistent with zero in
their analysis~C2---suggesting no residual isotropic
acceleration when the dipole is removed.  The defect
framework predicts $\Delta q(\beta, d) = (\beta/9)
(\lambda_H/d)^3 \approx 8.5$ at the WFH2009 primary amplitude (or
$\approx 6.6$ at the CF4++ amplitude), yielding
$\sim\!6\%$ (resp.\ $\sim\!18\%$) parameter-free agreement.
Colin et~al.~\cite{Colin2019} provided the precursor at
$3.9\sigma$ from the JLA sample, now confirmed at
$>5\sigma$.  Companion analysis:
Rameez~\cite{Rameez2024}.

The Tsagas (2026)~\cite{Tsagas2026ApJ} result that general
relativity doubles the peculiar-velocity growth rate compared
to Newtonian estimates amplifies the $\Delta H/H$ bridge prediction
by a factor $\sim 2$ at the SH0ES depth, strengthening the
pattern (though not the status, which remains EXPLORATORY).

Three issues prevent the tilted-FLRW programme from providing
a complete replacement for $\Lambda$CDM at this time.

\emph{(i) The $d^{-3}$ vs exponential form.}
The Tsagas power law $\Delta q \propto d^{-3}$ agrees with the
Colin exponential $e^{-z/S}$ within $20\%$ at
$d \approx 200$--$400\;\mathrm{Mpc}$ but diverges sharply
below $d \lesssim 150\;\mathrm{Mpc}$
, where the self-consistent
closure breaks down.  Redshift-binned analysis of the Rubin LSST
SN sample can distinguish the two functional forms, providing a
direct test of the scale-dependence prediction.

\emph{(ii) Quantitative $\vartheta$ calibration.}
The self-consistent closure $|\vartheta|/H = \beta\lambda_H/d$
assumes a single coherent monopolar flow.  The real velocity
field is filamentary, multi-scale, and direction-dependent.  A
quantitative comparison requires the redshift-weighted integral
of the Tsagas formula against the actual survey selection function
and the CF4 velocity-field decomposition.

\emph{(iii) BAO discrimination.}
At $d \gtrsim 500\;\mathrm{Mpc}$ (BAO scales), the tilt
correction drops below $0.5\%$ and all observers agree on $q$
to better than one part in $10^3$.  The BAO-measured expansion
history is therefore immune to the tilt mechanism.  If future BAO
analyses (DESI DR3+, Euclid spectroscopic) independently confirm
$w = -1$ at high precision, the tilt mechanism cannot substitute
for $\Lambda$---it can only modify the local interpretation of
the SN Hubble diagram.  If instead BAO data continue to prefer
$w_0 \neq -1$ (as in DESI DR2), the tilt + age-bias + evolving
$w$ decomposition~\eqref{eq:disc-q-decomp} becomes the natural
interpretive framework.

\emph{Data-honest synthesis.}
The Tsagas programme makes three qualitative predictions---dipolar
$q$, anomalous bulk flow, directional $H_0$ variation---all of
which have observational support at $3.9$--$5.7\sigma$.  The
defect framework contributes the algebraic ceiling on the tilt
amplitude ($\betasr \leq \eps_1/(1 + \etaud)$), the bridge
parameter $\etaud = 1/12$ connecting the two frameworks, and the
peculiar Jeans length defining the contamination boundary.  The
quantitative programme is limited by the self-consistent closure
(which overestimates at $d < 300\;\mathrm{Mpc}$; see
Section~\ref{sec:disc-backreaction} for the connection to
Buchert's tilted-flow averaging formalism), the unresolved
directional offset ($65^\circ$), and the model-dependence of the
Son et~al.\ BAO+CMB $q_0$.  These three limitations define the
frontier for future work.

% === END INLINED ===


% ============================================================
\section{Homogeneous framework and inhomogeneous reality}
\label{sec:disc-backreaction}
% ============================================================

The constraints derived in this work apply to the coherent,
spatially homogeneous component of cosmic anisotropy.  In the
real universe, structure formation generates kinematic
anisotropy that is spatially inhomogeneous, scale-dependent,
and epoch-dependent.


\subsection{The Buchert equations}
\label{sec:buchert-equations}

The scalar averaging of an irrotational dust cosmology over a
compact domain $\mathcal{D}$ yields effective Friedmann
equations with a kinematical backreaction
$\mathcal{Q}_\mathcal{D} = (2/3)
(\langle\theta^2\rangle_\mathcal{D}
- \langle\theta\rangle_\mathcal{D}^2)
- 2\langle\sigma^2\rangle_\mathcal{D}$~\cite{Buchert2000}.
The averaged shear
$\langle\sigma^2\rangle_\mathcal{D}$ is positive-definite and
never vanishes, even when the coherent shear tensor
$\langle\sigma_{ab}\rangle_\mathcal{D}$ averages to zero by
statistical isotropy.  The departure parameter $x$ constrains only
the coherent part; the variance part is invisible to the
Bianchi framework.


\subsection{Scale dependence}

The magnitude of the backreaction is critically
scale-dependent and foliation-dependent.  Three layers of
complexity must be distinguished: the mathematical
non-commutativity of averaging and evolution (a formal
property of general relativity), the foliation and frame
problem (which hypersurface and which congruence one
averages over), and the observational closure problem
(connecting averaged quantities to what observers actually
measure).

Numerical-relativity simulations on preferred hypersurfaces
find negligible \emph{box-average} backreaction, with
$\mathcal{Q}_\mathcal{D}/(6H^2) \sim 10^{-8}$ in
Poisson-like gauges~\cite{Aurrekoetxea2025}.
Oestreicher \& Koksbang~\cite{OestreicherKoksbang2024},
however, demonstrate that within $\sim\!200\,$Mpc sub-volumes
the effective curvature parameter $\Omega_R$ fluctuates by up
to $\sim\!10\%$, and voids expand $10$--$30\%$ faster than the
box average---a concrete realisation of the ``global small,
local large'' dichotomy.  In comoving synchronous gauge the
same simulations yield
$\mathcal{Q}_\mathcal{D}/(6H^2) \sim 10^{-1}$ on
$\sim 64\,h^{-1}\,$Mpc scales, a three-to-five order-of-magnitude
foliation dependence that reflects the genuine non-uniqueness
of spatial averaging in general relativity.

The Wiltshire timescape programme~\cite{Wiltshire2024}
provides explicit late-time magnitudes:
$|\bar{\Omega}_Q| \sim 10^{-6}$--$10^{-5}$ at last
scattering, growing to $\sim\!0.03$--$0.04$ by the present
epoch as void kinetic energy accumulates.  Pantheon+
reanalyses within this framework are competitive with
$\Lambda$CDM under some pipelines, though the Newtonian
$N$-body phenomenology (R\'acz AvERA: $H_0 \approx
71$--$73\;\mathrm{km\,s^{-1}\,Mpc^{-1}}$ in high-resolution
runs) does not constitute a full general-relativistic
derivation.

Buchert's scalar-averaging formalism has recently been extended
to general 3+1 foliations with tilted
flow~\cite{BuchertMourierRoy2020}---directly relevant to the
present framework, where the tilt rapidity $\beta$ defines
the frame mismatch between geometry and matter congruences.
A systematic connection between the departure parameter $x$ and the
Buchert backreaction
$\mathcal{Q}_\mathcal{D}$---perhaps through a coarse-graining
procedure that respects the tilt frame---remains an open
problem (Section~\ref{sec:open-problems}).

On sub-Hubble scales ($\lesssim 50\,$Mpc), local shear is of
order unity and the effective anisotropy vastly exceeds the
Bianchi constraint $x \sim 10^{-6}$.  On scales above the
statistical homogeneity threshold
($\sim 100$--$180\,$Mpc), structure-induced anisotropies
average out and the effective contribution drops below $x$.
The crossover scale at which backreaction-induced anisotropy
falls below the CMB-derived Bianchi bound is likely near or
slightly above the homogeneity scale.


\subsection{The Green--Wald no-go and its limitations}

The Green \& Wald~\cite{GreenWald2011,GreenWald2014}
perturbative argument shows that backreaction is small in the
weak-field limit.  Their framework uses a distributional weak
limit---a local operation well-defined at each spacetime
point---rather than the non-local spatial averaging that defines
backreaction in the Buchert programme.  The backreaction may be
inherently non-perturbative: even though the Newtonian potential
$\Phi \sim 10^{-5}$, density perturbations $\delta \gg 1$ on
sub-Mpc scales, and metric derivatives become large as structure
goes nonlinear.


\subsection{Practical assessment}

For the defect-framework scenarios: the Bianchi CMB bound at
$x \sim 10^{-6}$ is not contaminated by backreaction at
observationally relevant levels.  Global backreaction contributes
at most $\sim 10^{-5}$ to effective anisotropy in
observationally relevant gauges, comparable to but not exceeding
the Bianchi constraint.  The Bianchi bounds can therefore be
treated as genuine constraints on primordial anisotropy, with
backreaction corrections entering at a level below the current
observational sensitivity.

The honest acknowledgment is that the homogeneous Bianchi
framework provides a necessary but not sufficient description
of cosmic anisotropy.  Inhomogeneous anisotropy
modes---spatially varying tilt fields, local bulk flows,
effective anisotropy from super-cluster dynamics---are entirely
unconstrained by the bounds derived here.  We return to this
point in Section~\ref{sec:disc-limitations}.


% ============================================================
\section{Cosmic web vorticity and the scale hierarchy}
\label{sec:disc-cosmic-web}
% ============================================================

\subsection{The rotating filament}

Tudorache et~al.~\cite{Tudorache2025} report the detection of
a $\sim 15\,$Mpc rotating galaxy filament at $z = 0.032$,
comprising 14 HI-selected galaxies from the MIGHTEE-HI survey
with significant spin--filament alignment
($\langle|\cos\psi|\rangle = 0.64 \pm 0.05$, exceeding
simulation predictions).  The detection constitutes the third
HI-detected filamentary structure in the literature and the
largest individual rotating structure with spectroscopically
resolved kinematics.


\subsection{Perturbative versus background vorticity}

The filament rotation is perturbative, inhomogeneous vorticity
generated by structure formation---not background rotation in
the Bianchi sense.  The vorticity power spectrum has been
quantified in relativistic $N$-body simulations:
Jeli\'c-Cizmek et~al.~\cite{JelicCizmek2018} measure a peak
at $k_p \approx 2\,h\,$Mpc$^{-1}$ at $z = 1$, shifting to
$k_p \approx 1\,h\,$Mpc$^{-1}$ by $z = 0$ with peak amplitude
$P_\omega/(Hf)^2 \approx 5\;(\mathrm{Mpc}/h)^3$.
S\o rensen, Hannestad \& Tram~\cite{Sorensen2025} confirm these
results with an independent $N$-body code and demonstrate
estimator convergence to $\sim\!5\%$ accuracy across three
decades in $k$.  Wang et~al.~\cite{Wang2021} report
observational evidence for filament spin at
$\pm 100\;\mathrm{km\,s^{-1}}$ on Mpc scales with a
Hubble-timescale period, consistent with the
simulation-predicted vorticity amplitudes.

A rough estimate of the local vorticity from the Tudorache
filament gives $\omega_{\mathrm{local}}/H_0 \sim
v_{\mathrm{rot}}/r
\sim 10^5\;\mathrm{m\,s^{-1}}/(3 \times 10^{22}\;\mathrm{m})
\sim \mathcal{O}(1)$, exceeding the Saadeh background bound of
$(\omega/H)_0 < 5 \times 10^{-11}$ by ten orders of magnitude.

This discrepancy is not a contradiction.  The Saadeh bound
constrains the coherent, spatially uniform rotation of the
entire Bianchi background geometry, evaluated at
$z \sim 1100$ and projected onto the full CMB sky.  The
filament rotation and the vorticity power spectrum are
nonlinear, small-to-intermediate scale phenomena generated by
shell-crossing at expanding/collapsing boundaries---the
Weyl-magnetism mechanism in 1+3 covariant
language~\cite{Umeh2023}.  The Bianchi/CMB limits constrain
homogeneous large-scale modes; cosmic-web vorticity populates
an entirely different sector of the scale hierarchy.
These two regimes are not in tension.


\subsection{The scale hierarchy}
\label{rem:scale-hierarchy}

The complete hierarchy of vorticity scales, spanning thirteen
orders of magnitude, is:
\begin{center}\footnotesize
\begin{tabular}{ll}
\toprule
$\omega/H_0 \sim \mathcal{O}(1)$ &
  Tudorache filament ($15\,$Mpc) \\
$\omega/H_0 \sim 10^{-1}$ &
  Filament spin: $\pm 100\;\mathrm{km\,s^{-1}}$ at Mpc scale \\
$P_\omega/(Hf)^2 \approx 5\;(\mathrm{Mpc}/h)^3$ &
  Vorticity power spectrum peak ($k_p \approx 1\,h\,$Mpc$^{-1}$, $z = 0$) \\
$\omega/H_0 < 10^{-3}$ &
  MES algebraic bound (type-independent) \\
$(\omega/H)_0 < 5 \times 10^{-11}$ &
  Saadeh CMB bound (model-dependent) \\
$\omega/H_0 \sim 10^{-11}$ &
  Bianchi~VII$_h$ posterior \\
\bottomrule
\end{tabular}
\end{center}
The seven-order gap between the MES algebraic bound and the
Saadeh model-dependent bound quantifies the information gain
from morphological template fitting.  The ten-order gap between
the Saadeh bound and the Tudorache filament quantifies the
distinction between background and perturbative vorticity.
The vorticity power spectrum populates the intermediate regime,
confirming that the perturbative and background sectors are
well separated in scale space.
the conditioned legacy figure gallery in Appendix~\ref{app:conditioned-legacy-figure-gallery} presents this hierarchy graphically,
spanning from the cosmic-web filament rotation at $\omega/H_0 \sim
\mathcal{O}(1)$ down to the Saadeh model-dependent bound at
$\sim 10^{-11}$.


\subsection{Implications for the framework}

Three implications are relevant.

First, the Tudorache filament provides an observational
calibration point for the backreaction limitation: effective
anisotropy on sub-Hubble scales vastly exceeds the Bianchi
bound, confirming that the crossover scale is at or above the
statistical homogeneity threshold.

Second, the spin--filament alignment exceeding simulation
predictions is qualitatively consistent with the enhanced
angular momentum transfer predicted by the Tsagas relativistic
velocity growth law $v \propto a^2$.  In the 1+3 covariant
language, the chain $E_{ab}$ (tidal field) $\to$ $\beta$ (tilt)
$\to$ $\hat{\omega}_{ab}$ (matter-frame vorticity) is the
covariant analogue of tidal torque theory.  We stress that this
connection is speculative (level: ``speculation'' in the VT-08
hierarchy) and would require dedicated relativistic simulations
to test quantitatively.

Third, the HI surveys (MIGHTEE, WALLABY, full SKA) that can
detect rotating filaments open a complementary observational
window.  The vorticity power spectrum is now predicted to
$\sim\!5\%$ accuracy across three decades in
$k$~\cite{Sorensen2025}; observational mapping from filament
catalogues, combined with the background bounds from CMB
polarisation (LiteBIRD), would enable a two-scale
characterisation of cosmic angular momentum that
independently tests the Weyl-magnetism generation mechanism.


% ============================================================
\section{Methodological limitations as design choices}
\label{sec:disc-limitations}
% ============================================================

Before closing this chapter, we state five methodological
limitations of the present framework.  Each is a deliberate
design choice, not an oversight, and we explain the reasoning
that led to it.


\subsection{Summary statistics versus pixel-level analysis}

The Bayesian evidence comparison operates at the
summary-statistic level: it uses $D_2$, $D_3$, $\eps_1$, and
$\beta$ as scalar inputs, discarding the full CMB sky map and
the directional information of the matter dipole.  The
Saadeh~et~al.~\cite{Saadeh2016b} analysis demonstrates that
pixel-level template fitting yields constraints up to six orders
of magnitude tighter for the shear and vorticity.

We adopted summary statistics because the present framework's
novelty lies in the multi-probe combination---incorporating
matter dipole and bulk-flow data that pixel-level analyses
cannot access---rather than in the morphological detail of the
CMB sky.  The two approaches are complementary
(Section~\ref{sec:saadeh-comp}), and a future joint analysis
combining pixel-level CMB templates with matter dipole
likelihoods would be the natural next step.  The directional
information is the principal casualty of this design choice:
the $\ln\mathcal{B} \approx +26.3$ is a direction-marginalised
Bayes factor that could strengthen or weaken under directional
constraints.


\subsection{What the data constrain and what they do not}
\label{sec:disc-constraints}

The $K$-identifiability audit
(Proposition~\ref{prop:K-audit} of Chapter~\ref{ch:framework})
provides a sharp answer to the question ``what does the current
analysis establish?''

\emph{Established (C1--C2).}\quad
The data constrain the tilt rapidity
$\beta = (1.37 \pm 0.12) \times 10^{-3}$, conditional on the
matter-dipole anomaly being physical.  The departure is
tilt-dominated: $\Otilt$ contributes $>99\%$ of $x$; shear
contributes $<10^{-12}$.  Non-zero departure ($\ln\mathcal{B}
= +26.4$) is robust to LOCO ablation, prior variation, and
correlation sweeps.

\emph{Conditional (C3).}\quad
Equivalence-class preference (Tflat over O1D) is conditional on
anomaly physicality: if $\geq 50\%$ of the CatWISE/Radio dipole
excess is survey artifact, $\ln\mathcal{B} \to \leq 0$.  The
filling fraction $\FF \approx 9\%$ varies by $7\times$ with the
assumed intrinsic dipole amplitude $\varepsilon_1^{\rm int}$
(Section~\ref{sec:scenario-tables} of
Chapter~\ref{ch:results}); it is a
\emph{linear-budget-normalised score}, not a proven occupancy.

\emph{Not established (C4).}\quad
Bianchi type identification is not established.  The
scalar-amplitude likelihood collapses 14 models into 7
equivalence classes
(Table~\ref{tab:class-collapsed}); intra-class
$\Delta\ln\mathcal{B}$ is a prior-volume artifact.  The
geometry norm $\mathcal{G}_u$ is not $K$-identifiable under the
current compression (Proposition~\ref{prop:K-audit}(iv)).
Breaking this requires directional information---the Tier~2
programme of Chapter~\ref{ch:future}.

\emph{Informative vs.\ constraint channels.}\quad
Of seven active channels, two---(b) CatWISE+Radio and (c) CF4
bulk flow---provide $107\%$ of the evidence gain; channel~(a)
imposes a $-2.3$~nat F\&Q penalty; channels~(d), (e), (f),
(h) contribute $0.0$~nats of model discrimination
(the conditioned legacy figure gallery in Appendix~\ref{app:conditioned-legacy-figure-gallery}).  The evidence is a
two-measurement result dressed in a seven-channel architecture.

\emph{Three-layer reading.}\quad
We summarise the analysis by separating three logically
independent layers.
\emph{Layer~A} (theorem proximity):
$\FC \approx 6.3\%$ under the S3 scenario and flat comparator,
meaning the observed departure occupies $\sim\!6\%$ of the
MES-certified safe region.
\emph{Layer~B} (active-channel impact):
the standardised displacement norm $I_{\mathrm{act}} \approx 6.5$,
dominated by the CF4 bulk-flow channel ($61\%$ of $I_{\mathrm{act}}^2$)
and the CatWISE/Radio dipole channel ($28\%$), confirms that the
signal is large despite the small theorem occupancy.
\emph{Layer~C} (identifiability gate):
the geometry false-positive rate remains $22\%$ (threshold: $<10\%$),
so geometry claims are blocked
(Proposition~\ref{prop:K-audit}(iv)).
The combination---small $\FC$, large $I_{\mathrm{act}}$,
blocked $G_{\mathrm{id}}$---is the honest summary of the
current analysis level.


\subsection{Claim-tiered reading of the VER2 exports}
\label{sec:disc-ver2-claimtiers}

The VER2 export family sharpens this hierarchy rather than relaxing it.
Appendix~\ref{app:ver2-crosswalk} records five manifest-backed result
packs, but only two of them are promoted into the main text.
Result Pack~A and Result Pack~D are \ConditionalStatus\ surfaces:
Pack~A provides an observer-neutral morphology proxy and Pack~D provides
a predictive-residual certificate with an attached TSC overlay.
Neither one is allowed to stand in for directional template detection,
HTT posterior odds, or a truth certificate about the underlying
geometry.

The other three packs remain explicitly diagnostic.
Result Pack~B is a local-versus-global discrimination matrix and is not
a Bayes factor.
Result Pack~C is a descriptive $x/Q/F/G$ card whose filling language is
still downgraded to a proxy score until the propagation claim gates are
closed.
Result Pack~E is a validation-coverage matrix, but every executable
campaign in that registry still carries a \texttt{warn} status and
therefore functions as a no-claim gate rather than a validated science
promotion.
This separation matters because the headline evidence statements of
Chapter~\ref{ch:results} are still owned by the explicitly typed
BASS$\rightarrow$HTT inference path, not by the export layer that
packages those products for manuscript use.


\subsection{Calibrated transfer functions versus direct
  Boltzmann computation}

The shear-to-quadrupole conversion uses transfer functions
$T_2^{\mathrm{decay}} = 5{,}247.5\;\mu\mathrm{K}/(\sigma/H)$ and
$T_2^{\mathrm{grow}} = 1.05\;\mu\mathrm{K}/(\sigma/H)$ calibrated
against the BASS VER04 cross-validation
($D_2(\Sigstd=10^{-6}) = 17.53\;\mu\mathrm{K}^2$)
using the corrected formula $D_2^{\mathrm{shear}} = (6/2\pi)(T_2 \sigma_H)^2$
(no $T_0$ factor; see VER05 physics fix in Appendix).\footnote{%
  The AniCLASS interface specification exists in the codebase
  as an interface contract and stub; a live oracle connection
  requiring a full $\ell$-space Boltzmann solve is the
  planned next step.  The transfer-calibration flag
  (\texttt{w\_tr}) marks this gap explicitly.}
Direct integration of the Boltzmann hierarchy for each parameter
combination (as ABSolve and AniCLASS do) would yield more precise
$D_\ell(\boldsymbol{\theta})$ predictions.

We adopted the calibrated approach because the transfer function
is linear in the shear amplitude ($D_\ell \propto \sigma_H^2$)
to high precision, so a single calibration point suffices for
all amplitudes within the MES-allowed range.  The $x_h$
dependence, which is non-trivial, is captured by the
$f_2(x_h)$ and $f_3(x_h)$ interpolation tables.  The residual
calibration uncertainty is $\sim 10\%$ for the decay mode and
$\sim 20\%$ for the growing mode, propagated through the
evidence computation as an additional error term.  Direct
Boltzmann computation within the nested-sampling loop would
increase runtime by a factor of $\sim 10^3$, making it
computationally prohibitive for the 14-model comparison.


\subsection{Homogeneous Bianchi framework versus
  inhomogeneous reality}

The entire analysis assumes spatially homogeneous anisotropy
(the Bianchi classification) applied to a universe that is
observably inhomogeneous on sub-Hubble scales.  The backreaction
analysis of Section~\ref{sec:disc-backreaction} shows that the
Bianchi bound $x \sim 10^{-6}$ is not contaminated by
structure-formation backreaction at observationally relevant
levels, but the framework cannot constrain inhomogeneous
anisotropy modes---spatially varying tilt fields, local bulk
flows, or effective anisotropy from super-cluster dynamics.

We adopted the homogeneous framework because it is the simplest
extension of FLRW that admits the four kinematic degrees of
freedom (shear, vorticity, acceleration, tilt) simultaneously,
and because the MES bounds are the most assumption-explicit
constraints linking CMB observations to spacetime geometry.
The honest statement is that the framework constrains only the
coherent (monopole) component of cosmic anisotropy; the
variance component requires inhomogeneous methods.


\subsection{Single barotropic fluid versus multi-component}

The tilt--acceleration relation
$\dot{u}_a = [w/(3(1+w))]\tilde{D}_a\!\ln\rho$ and the
resulting $\etaud = 1/12$ hold strictly for a single
barotropic perfect fluid.  In a multi-component universe
(radiation + dust + dark energy), the four-acceleration retains
independent degrees of freedom not captured by the tilt
rapidity alone.  The safe-route bound
$\beta \leq \eps_1/(1 + \etaud)$ is therefore a
\emph{necessary} condition for consistency with the observed
dipole, but may not be sufficient in multi-fluid cosmologies.

We adopted the single-fluid approximation because it yields the
tightest (most constraining) bound on $\beta$, and because the
multi-fluid extension requires specifying the relative tilt
between radiation and matter---an additional free parameter not
constrained by current data.


\subsection{Static analysis versus dynamical evolution}

The departure variables $\Sigstd$, $\Wstd$, $\Otilt$ are
evaluated at the present epoch $z = 0$ and compared with CMB
data at $z \sim 1100$ through the MES bounds.  The dynamical
evolution between these epochs is captured only through the
transfer functions (for the shear--quadrupole conversion) and
the decay rates (Section~\ref{sec:teff-ode}), not through a
full time-dependent solution of the Einstein--Boltzmann system.

We adopted the static approach because the MES bounds are
epoch-independent (they constrain the ratio $\sigma/\Theta$ at
any time), and because the filling fraction $\FF$ is defined as
a ratio of present-epoch quantities.  The redshift evolution
$\FF(z)$ analysed in Section~\ref{sec:filling-z} uses
phenomenological $\beta(z)$ models rather than solutions of the
coupled DCP equations, which would require specifying initial
conditions at decoupling---a task deferred to future work.


% === INLINED: figures_ch08.tex ===
% Figures for Chapter 8
% Updated: TF-D06, 2026-03-12
% Added: fig_colin_beta, fig_tilted_H0_depth, fig_q_decomposition
% Replaced: fig_scale_hierarchy (lambda_J, d_S added)

% fig_q_decomposition moved to ch09


% === ADDED 2026-04-24: research-plan extensions (F90, F109) ===

\subsection{Sky-map overview of dipole and anomaly directions}
\label{sec:sky-map-overview}

Two complementary sky maps complete the directional discussion.
the conditioned legacy figure gallery in Appendix~\ref{app:conditioned-legacy-figure-gallery} shows the apex
directions of all five literature dipole estimators
(Planck~2018, CatWISE~2025, NVSS~2011, RACS~2023, CF4~2023) on a
single Galactic Mollweide projection, with a grey shaded
\emph{consensus area} that marks the region containing four of the
five apexes within $35^{\circ} \times 22^{\circ}$.  The NVSS apex sits
$\sim 90^{\circ}$ away from the consensus and is the only outlier;
the radio-dipole literature has long noted this offset and attributes
it to residual local-structure contamination at $z \lesssim 0.05$.
The convergence of the four remaining estimators across more than
six orders of magnitude in survey depth (from CMB last-scattering
to the $z \lesssim 0.05$ CF4 sample) is the key empirical input
to the framework.

Beyond the dipole apexes, the CMB literature catalogues a number of
preferred-direction anomalies whose mutual alignment has been the
subject of more than a decade of debate.
the conditioned legacy figure gallery in Appendix~\ref{app:conditioned-legacy-figure-gallery} consolidates seven such
directions on a single sky map: the Cold Spot, the Axis-of-Evil
(quadrupole-octupole alignment), the CatWISE photometric axis, the
CF4 bulk-flow axis, the Bianchi$_{VII_h}$ preferred axis, the
hemispherical-asymmetry axis, and the parity-odd power axis.  The
Axis-of-Evil and the Bianchi$_{VII_h}$ axis are within $\sim 5^{\circ}$
of each other---the great-circle joining the two is shown as a dashed
arc---suggesting a single small-$\ell$ preferred direction anchored
near $(l, b) \approx (240^{\circ}, 60^{\circ})$.  The Cold Spot and
hemispherical-asymmetry axes are decoupled from this cluster, in
agreement with the cone-overlap matrix of
the conditioned legacy figure gallery in Appendix~\ref{app:conditioned-legacy-figure-gallery}.

% === END INLINED ===
